\documentclass{article}
\usepackage{hyperref,amsmath}
\begin{document}
\section*{A published replay-recovery baseline for independent progress}
\paragraph{Decision being evaluated.}
An old keyed request has not returned a result to the caller. A later
failure can change which server holds the completion record or performs
recovery. New requests that do not use that service should continue,
without changing which old operation an eventual retry denotes.
This check examines the retry and acknowledgement layer. Radio,
actuator effects, lease service, and server migration are not run in the
fixture. No claim about the complete recovery contract follows.

\paragraph{Published baseline.}
Lee, Park, Kejriwal, Matsushita and Ousterhout,
\href{https://doi.org/10.1145/2815400.2815416}{SOSP 2015, RIFL},
Sections 3--4, separate request identity, atomic durable completion,
retry rendezvous, and reclamation. They explicitly require an exact
completion result alongside the mutation, and a retry to find the same
object-associated record after migration. Their client-wide cumulative
acknowledgement does not assert that every request above the prefix
is still unfinished. It does forbid reclaiming a prefix across a hole
under that acknowledgement interface. RIFL supplies an established
comparison for this project's earlier digital-store exact-retry results,
not for atomically committing physical motion with a journal entry.

\paragraph{Pinned tracker boundary.}
The pinned implementation uses a 512-position sequence window for
\texttt{newRpcId}; this is the executable constant, not the illustrative
128 mentioned in a source comment. For ordinary single-ID sequential
allocations, if the oldest unreturned ID is $b$, the last permitted
new ID before completing it is $b+511$. Higher successful completions
cannot raise the cumulative acknowledgement beyond $b-1$. With nothing
already assigned above $b$, at most 511 new allocations follow it
before a new call must prod the old request. If some above-$b$ IDs were
already assigned, they consume the same positions. The test's
\texttt{tryFinish} stops the test at the first prod; this observes the
blocking dependency without imposing a measured timeout or inferred
liveness failure in the actual transport. Returning the old result
unblocks allocation.

Two native tracker instances, when authority to use distinct live
client identities already exists, avoid that cross-request dependency.
The independent sequence advances even while the old request is
unreturned. The source test can issue and complete 2,048 independent
requests from the other tracker without prodding the old one. The
construction works for arbitrary finite numbers in the ideal
immediate-response model within numeric ID limits. It does not imply
an unbounded lease or every arbitrary workload. A lease-context
change costs an additional live identity, independent acknowledgement
state and any required renewal/establishment traffic.

\paragraph{Server-side check.}
The extracted \texttt{checkDuplicate} and \texttt{processAck} bodies
operate on serial fixture state. In one client context, replying to
later requests while the old result is missing costs retained results.
In the second context the same positive acknowledgement legitimately
reclaims only results belonging to that context. The old identity
still returns the stored result, or still reports a duplicate in
progress if no result was recorded. The latter is not permission to
execute again. Treating the greatest received result as an
acknowledgement prefix makes the old retry stale, even if the earlier
operation was not recorded as finished. Assigning the old request a
fresh client lease instead is recognized as a new key, not a retry.
These are local negative controls, not a proposal to change the
published interface. Reclamation of a result also cannot answer a
new exact audit that is outside the original retry contract.

\paragraph{Consequence for an integrated trial.}
A comparison against only a single stranded acknowledgement lane is
too weak if separate identities are authorized and affordable in the
candidate mission. When all old evidence travels with its designated
object and the new lane already exists, ordinary RIFL-style rendezvous
and lane separation explain this positive independent progress. If the
available roles forbid a second identity, that separation is an
additional resource, not a fair same-interface speedup. If a second
fault removes the rendezvous but the effect survives, neither a new
lane nor the old client's missing acknowledgement shows whether an
effect occurred. A supported alternate path to the same stable scope
or an actual atomic outcome record must still be constructed.
This pins the baseline and its missing source obligation without
claiming an architectural gain.

\paragraph{Replay scope.}
The client runs the pinned bodies and class with logging and test
callbacks stubbed. The server runs the two named bodies, with lease
validation, log freeing, record storage and locking supplied by a
serial harness. No durability, server migration, or failure timing
can be inferred from the stubs. There are 49 separated client controls,
seven shifted shared-window controls, one explicit 511-result limit,
and ten server controls.
\end{document}
